Um die gegebene Aussage zu zeigen, betrachten wir zunächst die Definition der Matrixnorm. Die Norm einer Matrix \( A \in M_n(\mathbb{C}) \) ist definiert als

\[
\|A\| = \sup_{x \in \mathbb{C}^n, \|x\| = 1} \|Ax\|.
\]

Das Quadrat der Norm ist dann gegeben durch

\[
\|A\|^2 = \left( \sup_{x \in \mathbb{C}^n, \|x\| = 1} \|Ax\| \right)^2 = \sup_{x \in \mathbb{C}^n, \|x\| = 1} \|Ax\|^2.
\]

Da \(\|Ax\|^2 = (Ax)^* (Ax) = x^* A^* A x\), können wir schreiben

\[
\|A\|^2 = \sup_{x \in \mathbb{C}^n, \|x\| = 1} x^* A^* A x.
\]

Die Matrix \(A^* A\) ist hermitesch, da \((A^* A)^* = A^* A\). Für eine hermitesche Matrix \(B\) ist bekannt, dass der Ausdruck \(x^* B x\) für \(\|x\| = 1\) durch den größten Eigenwert von \(B\) maximiert wird. Daher gilt

\[
\sup_{x \in \mathbb{C}^n, \|x\| = 1} x^* A^* A x = \lambda_{\max}(A^* A).
\]

Daraus folgt

\[
\|A\|^2 = \lambda_{\max}(A^* A).
\]

%CHECK: Die Schlussfolgerung wurde korrigiert, um den Fehler zu beheben, dass \(\|A^* A\|\) nicht gleich \(\lambda_{\max}(A^* A)\) ist.